\section{Limitations and conclusion}
\label{sec:limitations}
% Alias so the appendix's variance note, which cites \ref{sec:discussion},
% still resolves now that 06_discussion.tex is retired.
\label{sec:discussion}

\paragraph{The trainable projection head.}
Here $g$ is a trainable, never-pruned $\mathrm{Linear}(d,128)$:
$262{,}000$ parameters. With the teacher branch frozen, the optimiser can lower
PrC, FiC and SnC through $g$ alone: the minimisation admits solutions that never
touch $\theta$. Both branches share the encoder in ordinary contrastive
learning; freezing one inserts a free adapter. That bounds the claim, not the
measurement: accuracies came after a full fine-tune with $g$ discarded, and
$\Delta$ still isolates the objective. The defensible claim is better sparse
models, not preserved representations. The variant licensing the stronger claim --- one frozen head shared by student
and teachers --- lost at every measured sparsity (Table~\ref{tab:ablation}): the
clean interpretation is not the better configuration, and we chose the numbers.

\paragraph{What $\Delta$ actually contains.}
Three things separate the arms besides the objective, and we name them rather
than let them be inferred. The 25-epoch mask-frozen fine-tune runs on the BaCP
arm only. The pruning-phase learning rate differs ($0.01$ against $0.1$;
Section~\ref{sec:experiments:vgglr}). And BaCP's contrastive terms require two
independently augmented views, so its pruning phase runs the SimCLR
crop/colour/blur stack while I.P.\ runs crop-and-flip. Every $\Delta$ we report
carries all four together. The third matters most for the trend we emphasise:
recovery training and stronger augmentation both buy little on an intact
network and a great deal on a damaged one, which is the shape we attribute to
the objective. Separating them needs a baseline given the same fine-tune,
learning rate and augmentation, and we have not run it.

\paragraph{Budget, data and scope.}
Our 50 pruning plus 25 fine-tuning epochs at batch 512 are nearly an order of
magnitude fewer updates than 160 \citep{liu2021granet} or 250 \citep{li2025east}
epochs at batch 128, so nothing here is a state-of-the-art claim and the
absolute accuracies sit below the published literature. Results cover two
datasets, both CIFAR-scale and single-modality --- CIFAR-10 on five backbones,
CIFAR-100 on ResNet-34 --- so generality across resolution, dataset scale and
modality is untested, and in particular not CAP's own language setting: the
transfer we demonstrate runs one way. Three seeds bound the run-to-run spread
against a \result{cifar10.noise.floor.lo}--\result{cifar10.noise.floor.hi}
point floor \citep{bouthillier2021variance}; the paired tests in
Section~\ref{sec:results} are over cells that share models and data, so they
describe this grid rather than a population. We did not re-run RigL, GraNet,
GraSP or SynFlow, so whether objective- and support-level interventions compose
is open.

\paragraph{Conclusion.}
We studied CAP's contrastive pruning objective \citep{xu2022cap} --- theirs, not
ours --- on vision backbones, under three criteria, to 99.9\% sparsity, against
a baseline matched in schedule, criterion and pruning budget: a controlled measurement of what the
retraining objective is worth past the point where cross-entropy alone can
rebuild a pruned network.
